%!TEX root = ../../thesis.tex

\section{States and Observations}

For a given PUT, we model the state space as the joint state of process memory,
kernel state plus register file and other execution state necessary to PUT execution.
This constitutes a full description of the PUT and therefore suffices to fully predict future
behavior of the PUT given new inputs.
Although this state space formulation is straightforward, it is too large to serve as a
base for a fuzzing implementation.
In practice, this state is contained in the execution environment and even though it
is accessible to us, we will not use it directly.
Adhering to the approach of \gls{cgf}, the fuzzer instead observes coverage information $\mathcal{C}$.
In our case, this can either be block or edge coverage.
We therefore model the observation function $o_F:\mathcal{S}\rightarrow\mathcal{C}$ to map
true PUT state $\mathcal{S}_F$ to coverage feedback $\mathcal{C}$.
Later on, we will extend the observation function to include more information
depending on the concrete situation.

\section{Actions}

EfficientZero requires a discrete action space of tractable size, usually in the
order of tens to hundreds of actions.
In contrast to RL algorithms like Q-Learning that can deal well with probabilistic
transitions, EfficientZero requires actions to have a deterministic effect to allow proper planning.
This rules out the direct use of probabilistic mutation operators like random
bit flips, bit insertions or structural mutations like random splicing.
Designing a deterministic set of mutation operators is not easy since the effect
of an action must be compatible with the observation model.

Consider for example a seed $s$ that has been assigned a power of
$4\leftarrow\textsc{AssignEnergy}(s)$ and where a
mutation operator $m$ generated four new seeds $m(s,4)=\{s_0,s_1,s_2,s_3\}$.
\todo{explore reconciliation strategies}
that must be reconciled into a single observation since the mutation operator is modeled
as a single action, even if it generates multiple test cases.
\todo{explore reconciliation strategies}

\section{Rewards}

We will design our reward signal as a combination of a stateless reward $r_\alpha$
and a stateful reward $r_\beta$.
The stateless reward for time step $t$ is calculated from the observation $o_t$
while the stateful reward is allowed to integrate information from previous observations
$o_1,\dots,\o_t$.
The optimization objective is similar to that of other \gls{cgf} fuzzers.
In general, the stateless reward will be some function of the immediate coverage,
steering the learner towards high-coverage states.
The stateful reward will be some function of the coverage history, giving a measure
of \emph{marginal coverage}, i.e., newly attained coverage in comparison to the
previous seed, steering the learner towards high-coverage actions.
This is similar to the formulation~\textcite{bottinger_deep_2018} use.

\begin{align*}
  r_\alpha:\mathcal{C}&\rightarrow\mathbb{R}, \\
  r_\beta:\cup_{n=0}^{\infty}\mathcal{C}^n&\rightarrow\mathbb{R}, \\
  r&=r_\alpha+\omega r_\beta
\end{align*}


% TODO: Define the contract the agent and the fuzzer meet at. Ground this in
% efficientzero::abstraction (Environment / Action / Observation / StepResult):
%   - one environment step = one (or several) execution(s) of the emulated PUT,
%   - reset = restoring the target to a known state,
%   - observation = a fixed-width vector derived from block-hit counts,
%   - reward = a scalar derived from coverage.
% Present as a table mapping MDP component -> fuzzing artefact, and forward-reference
% the two instantiations.
%
% Then state, as REQUIREMENTS on the execution platform rather than as facts about
% MOGI, the four properties everything downstream leans on. Keeping them here as
% demands of the formulation is what separates the design from its realization ---
% \cref{chap:implementation} shows that MOGI satisfies each one, and the forward
% reference should be explicit at each:
%   P1  an exact and cheap reset, so the MDP is genuinely episodic (\cref{sec:rl-mdp})
%       and per-episode experimentation on a stateful target is tractable at all
%       (\cref{sec:platform-snapshots}),
%   P2  a fixed-width coverage observation derived from executed basic blocks, with
%       the aliasing that hashing into a bounded array implies
%       (\cref{sec:platform-coverage}),
%   P3  determinism: the same action sequence yields the same coverage vector on
%       every execution (\cref{sec:platform-determinism}),
%   P4  a throughput budget --- emulation is slow, which is what makes sample
%       efficiency the relevant currency (\cref{sec:platform-throughput}).
%
% Be explicit that P2 makes this a POMDP presented to the agent as an MDP: a hashed
% block-hit vector is a heavily aliased view of the guest's true state. State what
% that costs here, once, rather than twice in the two formulation sections.
